% TOP_PROBLEM: {"id": 3060, "title": "The additive basis conjecture", "category_id": 2, "category": {"id": 2, "name": "combinatorics", "display_name": "Combinatorics", "description": "Counting problems, graph theory, discrete structures.", "slug": "combinatorics", "order_index": 2, "created_at": "2026-07-31T15:26:25.671Z"}, "status": "open", "problem_number": "OPG-150", "set_id": 12}
% FIRST_POSTED: 2026-10-01
% ATTEMPT_STATUS: unresolved
% UnsolvedMath ID 3060; source catalogue code OPG-150
% !TeX program = lualatex
% GPT-6 Astra Ultra research attempt; independent partial work, 2026-10-01.
\documentclass[11pt,a4paper]{article}
\usepackage[margin=25mm]{geometry}
\usepackage{fontspec}
\setmainfont{lmroman10-regular.otf}[BoldFont=lmroman10-bold.otf,
 ItalicFont=lmroman10-italic.otf,BoldItalicFont=lmroman10-bolditalic.otf]
\usepackage{amsmath,amssymb,mathtools}
\usepackage{unicode-math}
\setmathfont{latinmodern-math.otf}
\AtBeginDocument{\let\setminus\smallsetminus}
\usepackage{xurl,hyperref}
\hypersetup{colorlinks=true,urlcolor=blue}
\setcounter{secnumdepth}{0}
\setlength{\parindent}{0pt}
\setlength{\parskip}{5pt}
\setlength{\emergencystretch}{3em}
\begin{document}
\section*{The additive basis conjecture}\label{3060}
{\small UnsolvedMath ID 3060 · OPG-150 · Combinatorics · OpenGarden}
\subsection{Definitions and mathematical statement}
For a prime $p$ and an integer $n\ge1$, let $V=\mathbb F_p^n$.
An additive basis here is a finite multiset $M$ of vectors for which every
$v\in V$ is the sum of a submultiset of $M$. Each occurrence can be used
at most once; equal vectors occurring several times remain separately
available. The empty sum is allowed and equals the zero vector.
This definition differs from an ordinary vector-space basis, where
arbitrary field coefficients are permitted.

The conjecture asks for a positive integer $c(p)$ depending only on $p$
such that, for every $n$ and every sequence of vector-space bases
$B_1,\ldots,B_{c(p)}$ of $V$, their multiset union is an additive basis.
Writing $B_i=(b_{i1},\ldots,b_{in})$, the assertion is
\[
 \left\{\sum_{i=1}^{c(p)}\sum_{j=1}^n
                 \epsilon_{ij}b_{ij}:\epsilon_{ij}\in\{0,1\}\right\}=V.
\]
The bases need not be distinct and may share vectors. All sums and the
coefficients $0,1$ are interpreted in $\mathbb F_p$.

The constant must be independent of dimension and independent of the
chosen bases. A bound which increases with $n$ would not settle this
question. The source also suggests the stronger candidate $c(p)=p$;
this is distinguished from the principal existence assertion, which
allows any finite value depending on $p$. Saying that the union
``contains'' an additive basis is equivalent to saying that the union
itself has the subset-sum property, since additional occurrences can
always be left unused.
\subsection{Short English statement}
For each prime, do a fixed number of arbitrary vector-space bases supply every vector as a sum using each available occurrence at most once?
\subsection{Research attempt: a sharp aligned case and the Fourier dimension loss}
\paragraph{Scope and literature check.}
GPT-6 Astra Ultra, 1 October 2026. Equal vector occurrences remain
separately available. The target is a constant independent of $n$.
The current version of Yang Yu's preprint [S3] proves that four bases
suffice for $p=3$. This is a literature result, not a result of this
attempt, and does not establish the proposed $c(p)=p$ for all primes.
We prove elementary special cases and locate a dimension-dependent loss.

\paragraph{1. A necessary lower bound and its aligned counterpart.}
Take $m$ copies of the standard basis. If $m<p-1$, every subset sum
has first coordinate in $\{0,1,\ldots,m\}$, and misses the residue
$p-1$. Hence any universal $c(p)$ must satisfy $c(p)\ge p-1$.

In dimension one, any $p-1$ nonzero vectors suffice, even when they
are different. To see this, start with $A_0=\{0\}$ and successively
put $A_j=A_{j-1}\cup(A_{j-1}+b_j)$. If $A_{j-1}$ is nonempty and
proper, it cannot satisfy $A_{j-1}+b_j=A_{j-1}$: translation by
$b_j\ne0$ is a cycle through all $p$ field elements. Therefore
$|A_j|\ge\min(p,|A_{j-1}|+1)$, and $A_{p-1}=\mathbb F_p$.

The same proof handles arbitrary dimension when every given basis
contains one nonzero multiple of each member of one fixed basis
$(e_1,\ldots,e_n)$. Apply the one-dimensional argument independently
to the $p-1$ occurrences in each coordinate direction, and take the
union of the selected occurrences. This establishes the optimal
constant $p-1$ for these aligned bases. For $p=2$ every single vector
space basis already has the desired subset-sum property.

\paragraph{2. Fourier inversion supplies a general but inadequate bound.}
Assume $p$ is odd. Select every occurrence in $m$ arbitrary bases
independently with probability $1/2$, and let $X$ be the resulting sum.
For $\xi\in\mathbb F_p^n$, let
$\chi_\xi(v)=\exp(2\pi i\langle\xi,v\rangle/p)$.
The product formula for the characteristic function is
\[
 \mathbb E\chi_\xi(X)=
 \prod_{i=1}^m\prod_{j=1}^n
 \frac{1+\exp(2\pi i\langle\xi,b_{ij}\rangle/p)}2.
\]
For nonzero $\xi$, at least one scalar product is nonzero in each
basis, since a basis spans the whole space. Every corresponding factor
has modulus at most $\rho=\cos(\pi/p)<1$; all other factors have
modulus at most one. Consequently
$|\mathbb E\chi_\xi(X)|\le\rho^m$.

Orthogonality of the characters, obtained by multiplying the geometric
sums in the $n$ coordinates, gives for every $v$
\[
 \Pr(X=v)=p^{-n}\sum_\xi
 \chi_\xi(-v)\mathbb E\chi_\xi(X)
 \ge p^{-n}\bigl(1-(p^n-1)\rho^m\bigr).
\]
Thus all vectors occur whenever
$m>\log(p^n-1)/(-\log\rho)$. This is a genuine finite bound for
every collection of bases, but grows linearly with dimension.

\paragraph{Verification and first remaining gap.}
The probability calculation is over labelled occurrences, so coincident
basis vectors cause no ambiguity. Fourier terms may be complex; the
lower bound uses their absolute values and the real-valued inversion
formula. The zero character contributes exactly one. The proof cannot
drop the factor $p^n-1$: it counts all nonzero characters after a
triangle inequality. A dimension-free conclusion would require stronger
joint information about these characters or a different method. The
aligned case avoids this loss because coordinates can be filled separately.

\subsection{Final status}
UNRESOLVED for general $p$ and arbitrary dimension. The aligned case
is proved with optimal $p-1$ bases, and an explicit dimension-dependent
general bound is derived. The known $p=3$ theorem is attributed to [S3].
Completion estimate: partial results only; no global percentage.

\subsection{Sources}
\begin{itemize}
\item[S1] ULAM.AI and the UnsolvedMath Contributors, supplied problems.json, record 3060 (OPG-150), OpenGarden collection. Catalogue identity and source transcription; CC BY 4.0.
\url{https://huggingface.co/datasets/ulamai/UnsolvedMath}.
\item[S2] Open Problem Garden, The additive basis conjecture. Original problem and discussion; the mathematical formulation below is expanded from the supplied source record.
\url{http://www.openproblemgarden.org/op/the_additive_basis_conjecture}.
\item[S3] Y. Yu, \emph{The Permanent Rank of a Matrix (Part Three): Note on the Additive Basis Conjecture}, arXiv:2510.01300, current version checked 1 October 2026.
\url{https://arxiv.org/abs/2510.01300}.
\end{itemize}
\end{document}
